\section{并行硬件与编程模型}
\label{sec:hw}

本节对应 CS149 的前半部分。我们沿着「如何把晶体管变成吞吐」这条线，依次讨论多核、SIMD、硬件多线程、
cache coherence 与内存一致性，再把这些概念映射到 GPU 的 SIMT 模型上。面试中这一层的问题看似基础，
但它们是回答「为什么这个 kernel 慢」的唯一依据。

\subsection{多核、SIMD 与硬件多线程}

现代处理器用三种正交的手段提高吞吐~\cite{hennessy2019golden}：（1）\textbf{多核}，复制整个取指--执行流水线；
（2）\textbf{SIMD}，一条指令驱动 $W$ 个 ALU（AVX2 为 8 个 fp32 lane，AVX-512 为 16 个），把取指译码开销摊到 $W$ 个数据上；
（3）\textbf{硬件多线程}（SMT），在一个核上保存多个线程的上下文，一个线程等待内存时切换到另一个线程。
三者分别利用了任务并行、数据并行与延迟隐藏。GPU 把三者推到极致：H100 有 132 个 SM，每个 SM 同时驻留最多 64 个 warp，
每个 warp 32 个线程以 SIMT 方式执行~\cite{lindholm2008tesla,nickolls2008cuda}。

\textbf{SIMD 发散}（divergence）是数据并行的主要效率损失。一组 $W$ 个 lane 必须执行到组内最慢的 lane 结束，因此利用率为
\begin{equation}
U = \frac{\sum_{i} w_i}{W\cdot \max_{i\in\text{gang}} w_i},
\end{equation}
其中 $w_i$ 为第 $i$ 个 lane 的有效迭代数。\lab{CS149 A1}的模拟显示，每个 8 宽 gang 中只要有一条慢 lane，
利用率就跌到 0.163（加速 1.30 倍而非 8 倍）；在 numpy 中实测，掩码执行在「每 16 个元素有 1 个慢元素」的输入上与全部都慢的输入
一样贵（353 对 356~ns/元素），先压缩活跃集合后降到 29~ns。GPU 的 warp 发散与此完全同构，工程解法是按工作量排序或分桶，
让同一个 warp 内线程的工作量接近，或者改为持续从队列拉取工作的 persistent kernel。

\begin{qa}{SIMD/warp 发散的代价怎么算？if/else 只执行一个分支，为什么还有代价？}
\ans 同一 warp 的线程共享一个程序计数器（Volta 之后有独立线程调度，但发散路径仍串行化），两个分支依次执行，
不活跃的 lane 被掩码。代价为 $W\cdot\max w_i/\sum w_i$ 倍。工作量差异越大，代价越高。
\wrong 「if/else 两个分支只执行一个，所以没有代价」；或者认为发散只发生在 GPU 而不发生在 CPU SIMD 上。
\end{qa}

\textbf{硬件多线程}用并发换延迟：若访存延迟为 $\ell$ 周期、每个线程两次访存之间平均有 $c$ 个周期的计算，
则约需 $\ell/c$ 个线程才能让核保持忙碌。代价是每个线程都要占用寄存器与上下文，片上存储被分摊。
GPU 的 occupancy（每个 SM 驻留 warp 数与上限之比）就是这件事的量化。

\subsection{Cache coherence、false sharing 与同步}

多核共享内存时，每个核私有的 cache 需要一致性协议。MESI 一类基于失效的协议~\cite{papamarcos1984low} 保证单个地址上
「写串行化」：一个核写某 cache line 前必须使其它核的副本失效。由此产生了 \textbf{false sharing}：
两个核写同一 64 字节 cache line 上的不同变量，协议让这条 line 在两个核之间来回迁移（每次数十到上百 ns），
表现为加线程反而变慢。修法是按 cache line 填充或让每个线程写私有副本，最后归约。

一致性（coherence）只约束单个地址；\textbf{内存一致性模型}（consistency）约束不同地址之间读写的可见顺序。
顺序一致性（SC）~\cite{lamport1979make} 要求所有核看到同一个全局交错顺序；x86 的 TSO 允许「写后读」重排（store buffer），
ARM 与 GPU 采用更弱的模型~\cite{adve1996shared}。程序员通过 fence 或 acquire/release 原子操作建立
happens-before 关系。GPU 上的典型陷阱是：用一个普通标志位在 block 之间传递「数据已就绪」，没有 \texttt{\_\_threadfence()}
时另一个 SM 可能先看到标志后看到数据。

无锁数据结构依赖 CAS（compare-and-swap），其经典问题是 \textbf{ABA}：线程读到栈顶 A 后被挂起，其他线程弹出 A、弹出 B、
再压回 A，CAS 仍然成功，但 A 的 next 指针已失效。解法是带版本号的指针、hazard pointer 或 epoch 回收。

\begin{qa}{什么是 false sharing？怎么发现，怎么修？}
\ans 不同核写同一 cache line 上的不同变量，coherence 协议使该 line 反复失效迁移。症状是线程数增加而吞吐下降，
perf 中 HITM（命中其它核修改过的 line）事件很高。修法：按 64~B 对齐填充每线程计数器，或写线程私有副本再归约。
\wrong 「没有共享变量就没有竞争」——竞争的单位是 cache line，不是变量。
\end{qa}

\begin{qa}{x86 上一个线程写 \texttt{data=1; flag=1}，另一个线程 \texttt{while(!flag); print(data)}，一定打印 1 吗？在 ARM 或 GPU 上呢？}
\ans x86 的 TSO 不重排「写--写」与「读--读」，所以硬件层面会打印 1，但编译器仍可能重排或把 \texttt{flag} 提到寄存器里，
必须用 \texttt{std::atomic} 的 release/acquire。ARM 与 GPU 的弱内存模型允许写--写重排，没有 fence 时可能打印 0。
\wrong 「有 cache coherence 就一定看到最新值」——coherence 不约束不同地址之间的顺序。
\end{qa}

同步原语的选择是另一个高频题。\lab{CS149 A2}在 Python 中复刻了 CS149 的任务系统，测得 8 个空任务一次 launch：
条件变量（sleep）线程池 38.7~\textmu s，每次新建线程 3169~\textmu s，纯自旋线程池 31--50~ms。后者在 C++ 中本应最快，
但在 CPython 中自旋线程持有 GIL，要等 5~ms 的切换间隔才让出，空闲时还烧满 1 个核。这个实验的价值在于说明
\emph{「spin 还是 sleep」取决于自旋者是否阻塞干活的线程}，而不是一个与平台无关的定论。
同一实验中，异步依赖图让 48 个 9 任务的 launch 从 259~ms 降到 152~ms，接近理想比值 $2/(9/8)=1.78$
（实测 1.71 倍），这是「填满尾部空闲」的直接证据，与 CUDA stream 和 CUDA Graph 的动机相同。

\subsection{GPU 的 SIMT 编程模型}

CUDA 把计算组织为 grid--block--thread 三级，block 内的线程可以通过 shared memory 通信并用 \texttt{\_\_syncthreads()} 同步，
block 之间默认独立~\cite{nickolls2008cuda}。硬件上 block 被调度到 SM，SM 以 32 线程的 warp 为单位发射指令。
H100 的存储层次见表~\ref{tab:gpumem}：寄存器与 shared memory 的带宽比 HBM 高一个数量级以上，延迟低一个数量级，
所以几乎所有 GPU 优化都是把数据「搬近、复用、搬远」三步~\cite{volkov2010better,cutlass}。

\begin{table}[!t]
\centering
\caption{H100 SXM 存储层次的数量级（附录表~\ref{tab:hw} 的子集）}
\label{tab:gpumem}
\footnotesize
\begin{tabularx}{\columnwidth}{@{}l l l Y@{}}
\toprule
层级 & 容量 & 带宽量级 & 延迟与说明 \\
\midrule
寄存器 & 256~KB/SM & — & 1 周期；寄存器分块是 GEMM 复用的最内层 \\
Shared/L1 & $\le$228~KB/SM & $\sim$30~TB/s（全卡） & 20--30 周期；软件管理，32 个 bank \\
L2 & 50~MB & $\sim$10~TB/s & $\sim$200 周期；K/V 重读多命中于此 \\
HBM3 & 80~GB & 3.35~TB/s & $\sim$500 周期；roofline 的默认 $B$ \\
主机内存 & 数百 GB & 32--64~GB/s & 微秒级；经 PCIe，offload 的瓶颈 \\
\bottomrule
\end{tabularx}
\end{table}

\textbf{访存合并}（coalescing）：一个 warp 的 32 个线程访问连续且对齐的地址时，硬件合并为少量 32~B/128~B 事务；
按列访问行主序矩阵（stride 等于行宽）时每个线程一个事务，有效带宽下降一个数量级。常用的修法是借 shared memory 做转置中转。

\textbf{Bank conflict}：shared memory 分为 32 个 bank、每 bank 4 字节宽。同一 warp 的线程访问同一 bank 的不同地址会被串行化。
典型例子是 $32\times32$ 的 float 矩阵按列读取，所有线程落在同一 bank；声明为 $32\times33$（填充一列）即可错开。
同一地址的广播不冲突。

\begin{qa}{什么是 memory coalescing？给一个不合并的例子。}
\ans 合并的单位是 warp，看的是 32 个线程地址的模式：连续对齐时一个 128~B 事务就能服务整个 warp。
行主序矩阵按列读、结构体数组（AoS）只读一个字段、随机 gather 都不合并。修法：改数据布局（SoA）、shared memory 转置、
或者让相邻线程处理相邻元素。
\wrong 「同一个 block 内的访问会自动合并」；或把 coalescing 与 shared memory bank conflict 混为一谈。
\end{qa}

\begin{qa}{Shared memory bank conflict 怎么产生、怎么避免？}
\ans 同一 warp 中多个线程访问同一 bank 的\emph{不同}地址时串行化。按列访问 $32\times32$ float 矩阵是 32 路冲突；
填充为 $32\times33$ 使第 $i$ 行第 $j$ 列落在 bank $(33i+j)\bmod 32$，冲突消失。Hopper 的 TMA 与 swizzle 布局也是为了这件事。
\wrong 「bank conflict 发生在 global memory」；「所有线程读同一地址也冲突」（那是广播）。
\end{qa}

\textbf{Occupancy} 是延迟隐藏的手段而不是目标。Volkov 的经典报告~\cite{volkov2010better} 指出，
依靠指令级并行（ILP）和寄存器分块，低 occupancy 的 kernel 可以更快：高性能 GEMM 常让每个线程计算 $8\times8$ 以上的输出块，
占用大量寄存器，occupancy 只有 25\% 左右，但数据复用显著更高。CUTLASS~\cite{cutlass} 的层次化分块
（threadblock tile → warp tile → MMA tile）把这一思想系统化。

\begin{qa}{occupancy 越高越好吗？}
\ans 不是。occupancy 只是隐藏延迟的一种方式，另一种是 ILP 与异步拷贝。寄存器分块很重的 GEMM 和 FlashAttention
故意降低 occupancy 换取复用。判断标准是吞吐：Nsight Compute 显示 stall 原因主要是「等待内存」时提高 occupancy 有用，
主要是「等待 tensor core」时没用。
\wrong 「把 occupancy 拉到 100\% 性能就最好」。
\end{qa}

\subsection{Tensor core 与 Hopper 的异步执行}

Volta 引入 tensor core，每个时钟完成一个小矩阵乘加~\cite{nvidiavolta2017}；Ampere 增加 BF16/TF32 与 2:4 结构化稀疏~\cite{nvidiaa100,mishra2021accelerating}；
Hopper 增加 FP8、TMA（Tensor Memory Accelerator，异步批量搬运张量块）、warpgroup 级的异步 MMA 指令（wgmma）
以及 thread block cluster 与分布式 shared memory~\cite{nvidiah100}。这些特性的共同方向是：
\emph{把搬运与计算交给不同的硬件单元异步执行}，让软件流水线（producer 负责搬运、consumer 负责计算）成为可能。
FlashAttention-3 的 warp specialization 正是建立在这些特性之上~\cite{shah2024flashattention3}。

\begin{keypts}
\begin{itemize}
\item 三种并行：多核（任务）、SIMD/SIMT（数据）、多线程（延迟隐藏）；发散的代价按组内最慢 lane 计。
\item coherence 管单地址，consistency 管多地址顺序；false sharing 的单位是 cache line；无锁结构注意 ABA。
\item GPU 三件事：合并访存（warp 地址模式）、bank conflict（shared memory）、occupancy（手段不是目标）。
\item Hopper 的方向是异步：TMA 负责搬、wgmma 负责算，软件用 warp specialization 把两者流水起来。
\end{itemize}
\end{keypts}

\begin{research}
\begin{itemize}
\item 如何自动选择 warp specialization 的 producer/consumer 划分与流水深度？这是一个带资源约束（寄存器、shared memory）的
调度问题，目前多由专家手写（CUTLASS、FlashAttention-3）。
\item 设计一个微基准，测量 H100 上 shared memory bank conflict、L2 命中与 TMA 异步拷贝对一个 tile 化 GEMM 的各自贡献；
关键是每次只改变一个因素，并给出与 roofline 上界的差距分解。
\end{itemize}
\end{research}
